% !TEX root = ../main.tex

\section{Kwantowy silik Stirlinga}
Szczególnym przypadkiem momentu pędu jest cząstka o spinie \(s = \inv{2}\). Układ taki ma dwa stany, które zapisujemy skrótowo:
\begin{align}
	\ket{\uparrow} &= \ket{\inv{2},\inv{2}}, &
	\ket{\downarrow} &= \ket{\inv{2},-\inv{2}}.
\end{align}
Spin taki mają na przykład elektron, proton, neutron, atom srebra i inne. Moment magnetyczny cząstki zależny jest od jej spinu \cite{griffiths_kwant}:
\begin{equation}
	\vu{\mu} = \gamma \vu{S},
\end{equation}
gdzie \(\gamma\) jest współczynnikiem żyromagnetycznym, a \(\vu{S}\) jest spinowym momentem pędu. W polu magnetycznym hamiltonian ma postać:
\begin{equation}\label{eq:def_hamiltonian_magnetyczny}
	\ham = - \vu{\mu} \cdot \vb{B} = - \gamma \vb{B} \cdot \vu{S},
\end{equation}
gdzie \(\vb{B}\) jest indukcją pola magnetycznego. Wybieramy układ współrzędnych, aby pole magnetyczne było ułożone wzdłóż osi \(z\) i równe \(\vb{B} = B \vu{z}\). Hamiltonian można zapisać:
\begin{equation}
	\ham = - \gamma B \hat{S}_z.
\end{equation}
Energia jest wartością własną hamiltonianu:
\begin{equation}
	\ham \ket{s,m_s} =
	- \gamma B \hat{S}_z \ket{s,m_s} =
	- \gamma B \hbar m_s \ket{s,m_s}.
\end{equation}
Dwa stany energetyczne są więc równe:
\begin{align}
	E_{\uparrow}   &= - \frac{\gamma B \hbar}{2}, &
	E_{\downarrow} &=   \frac{\gamma B \hbar}{2}.
\end{align}

W równowadze termodynamicznej:
\begin{gather}
	p_{m_s} = \inv{Z} \exp\pab{-\frac{E_{m_s}}{k_B T}}, \\
	Z = \exp\pab{\frac{\gamma B \hbar}{2 k_B T}} + \exp\pab{-\frac{\gamma B \hbar}{2 k_B T}} = 2\cosh\pab{\frac{\gamma B \hbar}{2 k_B T}}, \\
	U = p_{\uparrow}E_{\uparrow} + p_{\downarrow}E_{\downarrow} = - \frac{\hbar\omega}{2} \tgh\pab{\frac{\hbar\omega}{2 k_B T}}.
\end{gather}
